
\section{Conclusion}
\label{sec:conclusion}
This paper presents a novel longitudinal framework for quantifying the social cost of corporate data breaches by bridging the gap between institutional data exposure and the externalized costs faced by IDT victims.
To accomplish this, we performed a comprehensive 13-year harmonization of all six waves of the Identity Theft Supplement (2008-2021) to the National Crime Victimization Survey, pooling a sample of over 41,000 IDT victims. Our primary contribution is the development of a multi-dimensional social cost model that captures the full social cost of IDT, and the application of this model to quantify the social costs of mega-breach corporate events. Our model monetizes direct financial OOP losses, the opportunity cost of lost time, lawyer costs, and healthcare expenditures associated with the physical and emotional distress of IDT victimization. We document a significant shift in IDT patterns over the study period, noting a transition from traditional financial account fraud toward the misuse of digital credentials. We also find that the social cost of IDT is decreasing over the study period (2008-2021). 

By pairing this victim data with a multi-stage augmented chronology of corporate data breaches (primarily from the PRC data breach chronology), we performed hypothesis testing to verify the link between massive data exposure and subsequent victimization. Our Wilcoxon signed-rank analysis confirms a statistically significant increase in IDT discovery when accounting for a discovery lag of 1-2 months. Furthermore, our longitudinal ``breach-to-victim'' conversion model reveals that while the volume of exposed records is increasing, the marginal utility of a single compromised record has decreased as the data market becomes more saturated. 

The application of our model to landmark security incidents reveals a failure in current regulatory penalties to capture the full scope of damage externalized to consumers. For the 2009 Heartland and 2013 Target breaches, the social cost (even when measured through a conservative short-term lens) exceeded institutional settlements by factors of 5 and 18, respectively. While the costs calculated for the 2017 Equifax breach suggest a narrowing gap between corporate settlements and the short-term social cost, our long-term projected social cost indicates that the social cost of such massive exposure could remain in the billions, with an estimated upper-bound cost of \$1.72 billion that more than doubles the corporate settlement. Ultimately, these results suggest that as the data market reaches saturation, the individual risk per record may decline but the total social burden remains a systemic crisis.


% \rev{To accomplish this, we conducted the}
% first \rev{(to our knowledge) full} longitudinal analysis to harmonize and pool all six waves of the Identity Theft Supplement (2008-2021), and quantify the changing social cost of IDT. 

% Our findings reveal that while more educated, higher income, and White individuals experience higher rates of victimization, a greater social burden is taken on by lower socio-economic status and minority populations, who bear greater unrecoverable losses, more lost time, and more health-related distress related to IDT.

% We document the shift in IDT types, theft methods, and discovery methods over the years, and note a transition from traditional financial fraud to the misuse of social media and email credentials, and other digital thefts. Finally, our integration of PRC breach data indicates that while the number of exposed records in data breaches is increasing, the conversion rate from breach to successful IDT  has generally remained stable with a slight decreasing trend over the years. This suggests that while the market of stolen records is growing more saturated, the value of unique records remains high. Despite this, our hypothesis testing reveals that there is a significant rise in IDT reports following a mega-breach event. When applied to real-world cases, we found that the 2013 Target breach generated costs 13 times higher than its settlement, the 2017 Equifax breach settlement effectively covered its estimated social cost. This indicates while the total social cost problem remains in the multi-billions, the gap between corporate penalties and the realized consumer harm may be decreasing as the market for stolen data becomes more saturated.  The large scale of modern day data breaches ensures that IDT remains a multi-billion dollar social problem that requires ongoing institutional and policy intervention. 



